\section{拐点位置关于常规隔离率的两种变化方向}
\label{app:lambda}
为验证式\eqref{eq:lambda-q0-sign}确实可正可负，取
\[
N=10^6,\quad \beta=\frac16,\quad q_0=\frac1{10},\quad
\gamma=\frac1{10},\quad c_0=\frac{20}{3},\quad I_0=100.
\]
此时 $\beta_2^0/\beta_1^0=3/5,S_c=N/10,\bar S=3N/40$。固定 $S^*=N/5$，对 $r\in\{21/20,2\}$ 分别构造
\[
S_0=rS^*,\qquad
\eta=I_0+\frac{\beta_2^0}{\beta_1^0}(S_0-S^*)+\rho_1\ln(S^*/S_0).
\]
两组参数均有 $S_0+I_0<N$、$I_0<\eta<I_{\max}^{no}$ 和 $S_c<2\bar S<S^*$。余下初始人数可置于不参与社区传播的仓室。在每个构造出的参数点求偏导时，固定该点的 $S_0,\eta$；以上关系仅用于构造参数点，不作为求导时的额外约束。

代入式\eqref{eq:lambda-q0-sign}，得到
\begin{equation}
\frac{\partial\lambda}{\partial q_0}
=\frac{40}{9(\ln5)^2}
\left[\ln\frac53-\ln3\left(\frac{16}{5}(r-1)-\frac65\ln r\right)\right].
\end{equation}
当 $r=21/20$ 时，圆括号中的量在 $0$ 与 $4/25$ 之间。结合 $\ln(5/3)>2/5$、$\ln3<2$，方括号大于 $2/25>0$。当 $r=2$ 时，圆括号中的量大于 $2$；由 $\ln3>2/3$、$\ln(5/3)<2/3$，方括号严格为负。这些估计均可由 $(x-1)/x<\ln x<x-1$（$x>1$）得到，故两种偏导数符号均在允许参数域中实现。

